\newpage

\section{\texorpdfstring{\textbf{Relatório Técnico
(Gerencial)}}{Relatório Técnico (Gerencial)}}\label{relatuxf3rio-tuxe9cnico-gerencial}

\textbf{Entregável:} Diagnóstico inicial (Fase 1 --- Dados de
Produção)\\
\textbf{Pessoa de referência:} \textbf{Gerente de Produção}\\
\textbf{Terceira pessoa jurídica:} \textbf{Maringá Ferro Ligas}

\newpage

\subsection{\texorpdfstring{\textbf{1. Contexto, objetivo e leitura
recomendada}}{1. Contexto, objetivo e leitura recomendada}}\label{contexto-objetivo-e-leitura-recomendada}

Este relatório consolida, em linguagem técnica voltada ao público
gerencial, o diagnóstico de \textbf{disponibilidade} e
\textbf{qualidade} dos dados atualmente disponíveis no acervo de
produção da \textbf{Maringá Ferro Ligas}. O objetivo é reduzir risco de
decisão e de interpretação ao preparar dados para análises e modelagem,
evitando que iniciativas de melhoria de processo e automação sejam
guiadas por dados incompletos, ambíguos ou estruturalmente
inconsistentes.

Para o \textbf{Gerente de Produção}, o relatório responde a quatro
perguntas práticas:

\begin{enumerate}
\def\labelenumi{\arabic{enumi})}
\tightlist
\item
  \textbf{Quais dados existem}, por domínio e por forno, e qual é a
  cobertura temporal observada;\\
\item
  \textbf{Qual o nível de qualidade} observado, usando critérios
  reconhecidos (DAMA);\\
\item
  \textbf{O que pode ser utilizado agora} e \textbf{o que requer
  correção} antes de uso analítico;\\
\item
  \textbf{Quais decisões} de padronização, governança e instrumentação
  elevam rapidamente a confiabilidade do dado.
\end{enumerate}

\newpage

\subsection{\texorpdfstring{\textbf{2. Escopo do acervo
analisado}}{2. Escopo do acervo analisado}}\label{escopo-do-acervo-analisado}

O diagnóstico considerou \textbf{124 arquivos}, organizados por
\textbf{domínio} e, quando aplicável, por \textbf{forno} (\textbf{F1} a
\textbf{F5}). A tabela a seguir resume volume e cobertura temporal
observados por domínio.

\begin{longtable}[]{@{}
  >{\raggedright\arraybackslash}p{(\columnwidth - 8\tabcolsep) * \real{0.1905}}
  >{\raggedleft\arraybackslash}p{(\columnwidth - 8\tabcolsep) * \real{0.1905}}
  >{\raggedleft\arraybackslash}p{(\columnwidth - 8\tabcolsep) * \real{0.1905}}
  >{\raggedleft\arraybackslash}p{(\columnwidth - 8\tabcolsep) * \real{0.1905}}
  >{\centering\arraybackslash}p{(\columnwidth - 8\tabcolsep) * \real{0.2381}}@{}}
\toprule\noalign{}
\begin{minipage}[b]{\linewidth}\raggedright
\textbf{Domínio}
\end{minipage} & \begin{minipage}[b]{\linewidth}\raggedleft
\textbf{Arquivos (n)}
\end{minipage} & \begin{minipage}[b]{\linewidth}\raggedleft
\textbf{Linhas totais (aprox.)}
\end{minipage} & \begin{minipage}[b]{\linewidth}\raggedleft
\textbf{Colunas médias}
\end{minipage} & \begin{minipage}[b]{\linewidth}\centering
\textbf{Período coberto}
\end{minipage} \\
\midrule\noalign{}
\endhead
\bottomrule\noalign{}
\endlastfoot
Consumo Fornos & 39 & 38.123.141 & 27.4 & 2018-01-01 a 2025-04-28 \\
Corridas & 40 & 114.949 & 34.1 & 2018-01-01 a 2025-04-29 \\
Eletrodo & 1 & 761 & 5.0 & 2021-01-02 a 2025-12-03 \\
Informações Diária & 40 & 14.232 & 47.2 & 2018-01-01 a 2025-04-29 \\
Supervisorio Forno 4 & 2 & 486.250 & 228.0 & 2023-09-29 a 2025-01-02 \\
Supervisorio Forno 5 & 2 & 242.099 & 25.5 & 2021-01-06 a 2024-12-05 \\
\end{longtable}

\textbf{Observação operacional:} os fornos \textbf{F1} e \textbf{F2} são
reconhecidamente mais antigos, com defasagem tecnológica em relação a
\textbf{F4} e \textbf{F5}. Isso se reflete na variação de estrutura,
granularidade e padronização --- inclusive dentro do mesmo domínio.

\newpage

\subsection{\texorpdfstring{\textbf{3. Metodologia: critérios de
qualidade adotados (DAMA) e como foram
medidos}}{3. Metodologia: critérios de qualidade adotados (DAMA) e como foram medidos}}\label{metodologia-crituxe9rios-de-qualidade-adotados-dama-e-como-foram-medidos}

A avaliação seguiu seis dimensões de qualidade de dados, usuais em
governança:

\begin{itemize}
\tightlist
\item
  \textbf{Completude:} proporção de campos efetivamente preenchidos.\\
\item
  \textbf{Validade:} proporção de valores parseáveis e compatíveis com
  regras mínimas de domínio (tipos, datas e ranges plausíveis).\\
\item
  \textbf{Acurácia (proxy):} ausência de referência metrológica direta
  para todos os sinais; adotou-se uma \textbf{proxy de plausibilidade
  estatística} baseada em taxa de outliers/anomalias. \textbf{Nesta
  proxy, menor é melhor}.\\
\item
  \textbf{Tempestividade:} presença e consistência de coluna temporal
  utilizável (timestamp/data) e compatibilidade com a granularidade do
  domínio. Quando a coluna temporal é ausente/ambígua, o resultado pode
  aparecer como \textbf{n/d}.\\
\item
  \textbf{Unicidade:} ausência de duplicidades conforme regra do
  domínio.\\
\item
  \textbf{Consistência:} estabilidade do schema (nomes/quantidade de
  colunas) e coerência interna entre arquivos do mesmo domínio.
\end{itemize}

\newpage

\subsection{\texorpdfstring{\textbf{4. Interpretação gerencial: o que os
dados permitem (e o que ainda
impedem)}}{4. Interpretação gerencial: o que os dados permitem (e o que ainda impedem)}}\label{interpretauxe7uxe3o-gerencial-o-que-os-dados-permitem-e-o-que-ainda-impedem}

O diagnóstico evidencia que, em vários domínios, \textbf{volume de dados
não é o gargalo}. O principal limitador é a \textbf{qualidade estrutural
e temporal} (padronização de colunas, tipos e timestamp), que aumenta
risco de interpretação e dificulta análises temporais confiáveis.

De forma gerencial, o uso prático se organiza em três níveis:

\textbf{Nível 1 --- utilizável com baixo risco (após ajustes pontuais):}
bases com validade/consistência altas e timestamp padronizável,
adequadas para séries temporais e correlações operacionais.

\textbf{Nível 2 --- utilizável com risco moderado (requer
padronização):} bases com boa estrutura geral, mas com lacunas,
divergências de tipos/unidades e/ou necessidade de formalizar chaves e
coluna temporal.

\textbf{Nível 3 --- alto risco (exige reprocessamento antes de uso):}
bases com problemas de parsing, estrutura inválida, duplicidades ou
completude extremamente baixa; o uso direto tende a produzir conclusões
erradas com ``confiança estatística'' artificial.

\newpage

\subsection{\texorpdfstring{\textbf{5. Quantidade de dados: mínimos
viáveis e ``ótimos'' como ordens de
grandeza}}{5. Quantidade de dados: mínimos viáveis e ``ótimos'' como ordens de grandeza}}\label{quantidade-de-dados-muxednimos-viuxe1veis-e-uxf3timos-como-ordens-de-grandeza}

\begin{small}
\setlength{\tabcolsep}{2pt}
\renewcommand{\arraystretch}{1.20}
\begin{tabularx}{\linewidth}{>{\RaggedRight\arraybackslash}p{3.2cm} >{\RaggedRight\arraybackslash}p{4.3cm} >{\RaggedRight\arraybackslash}p{4.3cm} >{\RaggedRight\arraybackslash}p{4.3cm}}
\toprule
Tipo de modelo & Escopo & Mínimo viável & Nível ótimo \\
\midrule
\makecell[l]{Physic-based (balanços)\\} & \makecell[l]{Corridas + Informações Diária\\} & \makecell[l]{6–12 meses contínuos com timestamp padronizado\\} & \makecell[l]{18–24 meses com calibração documentada\\} \\
\hline
\makecell[l]{Supervisionado (regressão/class.)\\} & \makecell[l]{Corridas + Info Diária + Consumo Fornos\\} & \makecell[l]{12 meses rotulados e limpos por forno, com chave temporal\\} & \makecell[l]{24–36 meses com labels confiáveis e drift monitorado\\} \\
\hline
\makecell[l]{Não supervisionado (detecção de padrão)\\} & \makecell[l]{Supervisório (F4/F5); senão Consumo Fornos\\} & \makecell[l]{2–3 meses contínuos com timestamp confiável\\} & \makecell[l]{6–12 meses com tags harmonizadas e filtros de ruído\\} \\
\hline
\makecell[l]{Reforço (RL)\\} & \makecell[l]{Supervisório com ações e retorno (F4/F5)\\} & \makecell[l]{3–6 meses de episódios com ação/recompensa registrada\\} & \makecell[l]{12+ meses com atrasos medidos e segurança validada\\} \\
\bottomrule
\end{tabularx}
\end{small}

\clearpage

\subsection{\texorpdfstring{\textbf{6. Quadro consolidado: qualidade por
forno e por domínio (uma
página)}}{6. Quadro consolidado: qualidade por forno e por domínio (uma página)}}\label{quadro-consolidado-qualidade-por-forno-e-por-domuxednio-uma-puxe1gina}

\begin{scriptsize}
\setlength{\tabcolsep}{1.4pt}
\renewcommand{\arraystretch}{1.05}
\begin{tabularx}{\linewidth}{>{\RaggedRight\arraybackslash}p{2.6cm} >{\RaggedRight\arraybackslash}p{1.0cm} >{\RaggedRight\arraybackslash}p{0.9cm} >{\RaggedRight\arraybackslash}p{1.0cm} >{\RaggedRight\arraybackslash}p{1.0cm} >{\RaggedRight\arraybackslash}p{1.0cm} >{\RaggedRight\arraybackslash}p{1.0cm} >{\RaggedRight\arraybackslash}p{1.0cm} >{\RaggedRight\arraybackslash}p{1.0cm} >{\RaggedRight\arraybackslash}p{1.4cm} X X}
\toprule
Domínio & Forno & n & Comp. & Val. & AccP & Temp. & Uniq. & Cons. & Coluna tempo & Interpretação & Ações imediatas \\
\midrule
\makecell[l]{Consumo Fornos\\} & \makecell[l]{F1\\} & \makecell[l]{7\\} & \makecell[l]{0.008\\} & \makecell[l]{0.824\\} & \makecell[l]{0.048\\} & \makecell[l]{n/d\\} & \makecell[l]{0.188\\} & \makecell[l]{0.714\\} & \makecell[l]{Data\\} & \makecell[l]{Risco alto: completude muito baixa\\Timestamp n/d; usar só após saneamento} & \makecell[l]{Sanear schema, deduplicar exportações\\Definir coluna temporal antes de modelar} \
\makecell[l]{Consumo Fornos\\} & \makecell[l]{F2\\} & \makecell[l]{8\\} & \makecell[l]{0.008\\} & \makecell[l]{0.824\\} & \makecell[l]{0.042\\} & \makecell[l]{n/d\\} & \makecell[l]{0.186\\} & \makecell[l]{0.625\\} & \makecell[l]{Data\\} & \makecell[l]{Risco alto: completude muito baixa\\Timestamp n/d; usar só após saneamento} & \makecell[l]{Sanear schema, deduplicar exportações\\Definir coluna temporal antes de modelar} \
\makecell[l]{Consumo Fornos\\} & \makecell[l]{F3\\} & \makecell[l]{8\\} & \makecell[l]{0.006\\} & \makecell[l]{0.823\\} & \makecell[l]{0.052\\} & \makecell[l]{n/d\\} & \makecell[l]{0.179\\} & \makecell[l]{0.5\\} & \makecell[l]{Data\\} & \makecell[l]{Risco alto: completude muito baixa\\Timestamp n/d; usar só após saneamento} & \makecell[l]{Sanear schema, deduplicar exportações\\Definir coluna temporal antes de modelar} \
\makecell[l]{Consumo Fornos\\} & \makecell[l]{F4\\} & \makecell[l]{8\\} & \makecell[l]{0.012\\} & \makecell[l]{0.826\\} & \makecell[l]{0.063\\} & \makecell[l]{n/d\\} & \makecell[l]{0.181\\} & \makecell[l]{0.625\\} & \makecell[l]{Data\\} & \makecell[l]{Risco alto: completude muito baixa\\Timestamp n/d; usar só após saneamento} & \makecell[l]{Sanear schema, deduplicar exportações\\Definir coluna temporal antes de modelar} \
\makecell[l]{Consumo Fornos\\} & \makecell[l]{F5\\} & \makecell[l]{8\\} & \makecell[l]{0.012\\} & \makecell[l]{0.826\\} & \makecell[l]{0.064\\} & \makecell[l]{n/d\\} & \makecell[l]{0.182\\} & \makecell[l]{0.75\\} & \makecell[l]{Data\\} & \makecell[l]{Risco alto: completude muito baixa\\Timestamp n/d; usar só após saneamento} & \makecell[l]{Sanear schema, deduplicar exportações\\Definir coluna temporal antes de modelar} \
\makecell[l]{Corridas\\} & \makecell[l]{F1\\} & \makecell[l]{8\\} & \makecell[l]{0.993\\} & \makecell[l]{0.799\\} & \makecell[l]{0\\} & \makecell[l]{n/d\\} & \makecell[l]{1\\} & \makecell[l]{1\\} & \makecell[l]{Data_Base\\} & \makecell[l]{Estrutura consistente; validade moderada\\Timestamp n/d exige padronização} & \makecell[l]{Padronizar parse/unidades e coluna temporal\\Documentar schema e validar consistência} \
\makecell[l]{Corridas\\} & \makecell[l]{F2\\} & \makecell[l]{8\\} & \makecell[l]{0.979\\} & \makecell[l]{0.798\\} & \makecell[l]{0\\} & \makecell[l]{n/d\\} & \makecell[l]{1\\} & \makecell[l]{0.625\\} & \makecell[l]{Data_Base\\} & \makecell[l]{Estrutura consistente; validade moderada\\Timestamp n/d exige padronização} & \makecell[l]{Padronizar parse/unidades e coluna temporal\\Documentar schema e validar consistência} \
\makecell[l]{Corridas\\} & \makecell[l]{F3\\} & \makecell[l]{8\\} & \makecell[l]{0.995\\} & \makecell[l]{0.785\\} & \makecell[l]{0.001\\} & \makecell[l]{n/d\\} & \makecell[l]{1\\} & \makecell[l]{1\\} & \makecell[l]{Data_Base\\} & \makecell[l]{Estrutura consistente; validade moderada\\Timestamp n/d exige padronização} & \makecell[l]{Padronizar parse/unidades e coluna temporal\\Documentar schema e validar consistência} \
\makecell[l]{Corridas\\} & \makecell[l]{F4\\} & \makecell[l]{8\\} & \makecell[l]{0.993\\} & \makecell[l]{0.78\\} & \makecell[l]{0.001\\} & \makecell[l]{n/d\\} & \makecell[l]{1\\} & \makecell[l]{1\\} & \makecell[l]{Data_Base\\} & \makecell[l]{Estrutura consistente; validade moderada\\Timestamp n/d exige padronização} & \makecell[l]{Padronizar parse/unidades e coluna temporal\\Documentar schema e validar consistência} \
\makecell[l]{Corridas\\} & \makecell[l]{F5\\} & \makecell[l]{8\\} & \makecell[l]{0.992\\} & \makecell[l]{0.78\\} & \makecell[l]{0.001\\} & \makecell[l]{n/d\\} & \makecell[l]{1\\} & \makecell[l]{1\\} & \makecell[l]{Data_Base\\} & \makecell[l]{Estrutura consistente; validade moderada\\Timestamp n/d exige padronização} & \makecell[l]{Padronizar parse/unidades e coluna temporal\\Documentar schema e validar consistência} \
\makecell[l]{Eletrodo\\} & \makecell[l]{F1 a F5\\} & \makecell[l]{1\\} & \makecell[l]{1\\} & \makecell[l]{0.799\\} & \makecell[l]{n/d\\} & \makecell[l]{0.993\\} & \makecell[l]{0.986\\} & \makecell[l]{1\\} & \makecell[l]{Data medicao\\} & \makecell[l]{Validade moderada; proxy de acurácia indefinida\\Precisa timestamp e calibração simples} & \makecell[l]{Padronizar unidades e chave de medição\\Calibrar sinais críticos e registrar timestamp} \
\makecell[l]{Eletrodo\\} & \makecell[l]{F1 a F5\\} & \makecell[l]{1\\} & \makecell[l]{1\\} & \makecell[l]{0.8\\} & \makecell[l]{n/d\\} & \makecell[l]{1\\} & \makecell[l]{1\\} & \makecell[l]{1\\} & \makecell[l]{Data medicao\\} & \makecell[l]{Validade moderada; proxy de acurácia indefinida\\Precisa timestamp e calibração simples} & \makecell[l]{Padronizar unidades e chave de medição\\Calibrar sinais críticos e registrar timestamp} \
\makecell[l]{Eletrodo\\} & \makecell[l]{F1 a F5\\} & \makecell[l]{1\\} & \makecell[l]{1\\} & \makecell[l]{0.8\\} & \makecell[l]{n/d\\} & \makecell[l]{0.977\\} & \makecell[l]{1\\} & \makecell[l]{1\\} & \makecell[l]{Data medicao\\} & \makecell[l]{Validade moderada; proxy de acurácia indefinida\\Precisa timestamp e calibração simples} & \makecell[l]{Padronizar unidades e chave de medição\\Calibrar sinais críticos e registrar timestamp} \
\makecell[l]{Eletrodo\\} & \makecell[l]{F1 a F5\\} & \makecell[l]{1\\} & \makecell[l]{1\\} & \makecell[l]{0.799\\} & \makecell[l]{n/d\\} & \makecell[l]{0.995\\} & \makecell[l]{1\\} & \makecell[l]{1\\} & \makecell[l]{Data medicao\\} & \makecell[l]{Validade moderada; proxy de acurácia indefinida\\Precisa timestamp e calibração simples} & \makecell[l]{Padronizar unidades e chave de medição\\Calibrar sinais críticos e registrar timestamp} \
\makecell[l]{Eletrodo\\} & \makecell[l]{F1 a F5\\} & \makecell[l]{1\\} & \makecell[l]{1\\} & \makecell[l]{0.8\\} & \makecell[l]{n/d\\} & \makecell[l]{1\\} & \makecell[l]{1\\} & \makecell[l]{1\\} & \makecell[l]{Data medicao\\} & \makecell[l]{Validade moderada; proxy de acurácia indefinida\\Precisa timestamp e calibração simples} & \makecell[l]{Padronizar unidades e chave de medição\\Calibrar sinais críticos e registrar timestamp} \
\makecell[l]{Informações Diária\\} & \makecell[l]{F1\\} & \makecell[l]{8\\} & \makecell[l]{0.808\\} & \makecell[l]{0.942\\} & \makecell[l]{n/d\\} & \makecell[l]{n/d\\} & \makecell[l]{1\\} & \makecell[l]{0.75\\} & \makecell[l]{Data_base\\} & \makecell[l]{Completude/validade medianas; schema varia\\Timestamp incerto reduz rastreabilidade} & \makecell[l]{Harmonizar schema/tipos e coluna temporal\\Fechar lacunas obrigatórias antes de uso} \
\makecell[l]{Informações Diária\\} & \makecell[l]{F2\\} & \makecell[l]{8\\} & \makecell[l]{0.839\\} & \makecell[l]{0.943\\} & \makecell[l]{n/d\\} & \makecell[l]{n/d\\} & \makecell[l]{1\\} & \makecell[l]{1\\} & \makecell[l]{Data_base\\} & \makecell[l]{Completude/validade medianas; schema varia\\Timestamp incerto reduz rastreabilidade} & \makecell[l]{Harmonizar schema/tipos e coluna temporal\\Fechar lacunas obrigatórias antes de uso} \
\makecell[l]{Informações Diária\\} & \makecell[l]{F3\\} & \makecell[l]{8\\} & \makecell[l]{0.834\\} & \makecell[l]{0.925\\} & \makecell[l]{n/d\\} & \makecell[l]{0.998\\} & \makecell[l]{1\\} & \makecell[l]{0.5\\} & \makecell[l]{Data_base\\} & \makecell[l]{Completude/validade medianas; schema varia\\Timestamp incerto reduz rastreabilidade} & \makecell[l]{Harmonizar schema/tipos e coluna temporal\\Fechar lacunas obrigatórias antes de uso} \
\makecell[l]{Informações Diária\\} & \makecell[l]{F4\\} & \makecell[l]{8\\} & \makecell[l]{0.736\\} & \makecell[l]{0.914\\} & \makecell[l]{0.001\\} & \makecell[l]{n/d\\} & \makecell[l]{1\\} & \makecell[l]{0.5\\} & \makecell[l]{Data_base\\} & \makecell[l]{Completude/validade medianas; schema varia\\Timestamp incerto reduz rastreabilidade} & \makecell[l]{Harmonizar schema/tipos e coluna temporal\\Fechar lacunas obrigatórias antes de uso} \
\makecell[l]{Informações Diária\\} & \makecell[l]{F5\\} & \makecell[l]{8\\} & \makecell[l]{0.768\\} & \makecell[l]{0.916\\} & \makecell[l]{n/d\\} & \makecell[l]{n/d\\} & \makecell[l]{1\\} & \makecell[l]{0.625\\} & \makecell[l]{Data_base\\} & \makecell[l]{Completude/validade medianas; schema varia\\Timestamp incerto reduz rastreabilidade} & \makecell[l]{Harmonizar schema/tipos e coluna temporal\\Fechar lacunas obrigatórias antes de uso} \
\makecell[l]{Supervisorio Forno 4\\} & \makecell[l]{F4\\} & \makecell[l]{2\\} & \makecell[l]{0.909\\} & \makecell[l]{0.994\\} & \makecell[l]{0.03\\} & \makecell[l]{n/d\\} & \makecell[l]{1\\} & \makecell[l]{1\\} & \makecell[l]{Timestamp\\} & \makecell[l]{Dados ricos; tempestividade parcial\\Schema volumoso exige governança} & \makecell[l]{Padronizar timestamp e tags essenciais\\Filtrar outliers e documentar alarmes} \
\makecell[l]{Supervisorio Forno 5\\} & \makecell[l]{F5\\} & \makecell[l]{1\\} & \makecell[l]{0.553\\} & \makecell[l]{0.024\\} & \makecell[l]{0.000\\} & \makecell[l]{1\\} & \makecell[l]{1\\} & \makecell[l]{1\\} & \makecell[l]{Data\\} & \makecell[l]{Validade muito baixa; registros heterogêneos\\Uso direto arriscado sem reprocessar} & \makecell[l]{Reprocessar e padronizar parse/timestamp\\Validar ranges e bloquear valores inválidos} \
\makecell[l]{Supervisorio Forno 5\\} & \makecell[l]{F5\\} & \makecell[l]{1\\} & \makecell[l]{0.956\\} & \makecell[l]{0\\} & \makecell[l]{n/d\\} & \makecell[l]{n/d\\} & \makecell[l]{1\\} & \makecell[l]{1\\} & \makecell[l]{Data\\} & \makecell[l]{Validade muito baixa; registros heterogêneos\\Uso direto arriscado sem reprocessar} & \makecell[l]{Reprocessar e padronizar parse/timestamp\\Validar ranges e bloquear valores inválidos} \
\bottomrule
\end{tabularx}
\end{scriptsize}

\clearpage

\subsection{\texorpdfstring{\textbf{7. Resumo gerencial: viabilidade de
modelagem por forno (Nível de
Risco)}}{7. Resumo gerencial: viabilidade de modelagem por forno (Nível de Risco)}}\label{resumo-gerencial-viabilidade-de-modelagem-por-forno-nuxedvel-de-risco}

Regras de cálculo (determinísticas): - Classificação por domínio/forno:
\textbf{Nível 1} quando há timestamp utilizável e nenhum indicador
crítico (comp≥0.85, val≥0.90, uniq≥0.90, cons≥0.70). \textbf{Nível 2}
quando a base é estruturada mas falta padronização temporal ou há
lacunas moderadas. \textbf{Nível 3} quando completude é muito baixa
(\textless0.30), validade \textless0.60, timestamp ausente/ambíguo ou o
domínio requerido não existe. - Dependências por tipo de modelo:
\textbf{Physic-based} = Corridas + Informações Diária;
\textbf{Supervisionado} = Corridas + Informações Diária + Consumo
Fornos; \textbf{Não supervisionado} = Supervisório (se existir) senão
Consumo Fornos; \textbf{Reforço (RL)} = Supervisório obrigatório. -
Aplicação aos domínios: Consumo Fornos = N3; Corridas = N2; Informações
Diária = N2; Eletrodo = N2; Supervisório F4 = N2; Supervisório F5 = N3.
O nível por forno/modelo é o pior (máximo) entre os domínios requeridos.

\begin{longtable}[]{@{}
  >{\raggedright\arraybackslash}p{(\columnwidth - 8\tabcolsep) * \real{0.2000}}
  >{\raggedright\arraybackslash}p{(\columnwidth - 8\tabcolsep) * \real{0.2000}}
  >{\raggedright\arraybackslash}p{(\columnwidth - 8\tabcolsep) * \real{0.2000}}
  >{\raggedright\arraybackslash}p{(\columnwidth - 8\tabcolsep) * \real{0.2000}}
  >{\raggedright\arraybackslash}p{(\columnwidth - 8\tabcolsep) * \real{0.2000}}@{}}
\toprule\noalign{}
\begin{minipage}[b]{\linewidth}\raggedright
Forno
\end{minipage} & \begin{minipage}[b]{\linewidth}\raggedright
Physic-based
\end{minipage} & \begin{minipage}[b]{\linewidth}\raggedright
Supervisionado
\end{minipage} & \begin{minipage}[b]{\linewidth}\raggedright
Não supervisionado
\end{minipage} & \begin{minipage}[b]{\linewidth}\raggedright
Reforço (RL)
\end{minipage} \\
\midrule\noalign{}
\endhead
\bottomrule\noalign{}
\endlastfoot
F1 & Nível 2 & Nível 3 & Nível 3 & Nível 3 \\
F2 & Nível 2 & Nível 3 & Nível 3 & Nível 3 \\
F3 & Nível 2 & Nível 3 & Nível 3 & Nível 3 \\
F4 & Nível 2 & Nível 3 & Nível 2 & Nível 2 \\
F5 & Nível 2 & Nível 3 & Nível 3 & Nível 3 \\
\end{longtable}

Legenda: Nível 1 = pronto para uso; Nível 2 = uso com risco moderado
(padronização necessária); Nível 3 = alto risco ou ausência de domínio
requerido.

\subsection{\texorpdfstring{\textbf{8. Conclusões e decisões que o
Gerente de Produção precisa
tomar}}{8. Conclusões e decisões que o Gerente de Produção precisa tomar}}\label{conclusuxf5es-e-decisuxf5es-que-o-gerente-de-produuxe7uxe3o-precisa-tomar}

Este diagnóstico permite decisões objetivas sobre priorização, com
impacto direto na capacidade de fazer análises e sustentar modelos com
credibilidade operacional.

\textbf{Decisão 1 --- Prioridade de curto prazo (0--4 semanas):} definir
quais domínios serão tratados como ``base mínima'' para o próximo ciclo.

\textbf{Decisão 2 --- Padronização mandatória:} estabelecer padrão
corporativo simples para \textbf{timestamp}, \textbf{timezone}, nomes de
colunas e unidades. Sem isso, análises temporais ficam expostas a erro
de interpretação.

\textbf{Decisão 3 --- Correção estruturante (reprocessamento):} aprovar
a revisão de rotinas de extração onde a estrutura atual gera risco
elevado (parsing, completude e duplicidades).

\textbf{Decisão 4 --- Instrumentação do alvo (distância do eletrodo):}
como não há medição direta da variável de interesse no acervo, é
necessário decidir o caminho: \textbf{proxies (curto prazo)} ou
\textbf{medição/inferência instrumentada (médio prazo)}.

\newpage
